\exercises

\tierunderstand

\begin{bt}
  \item Phương trình~\eqref{eq:ch06-tong-co-trong-so} có chỉ số \(i\) ở vế
    trái, còn phương trình~\eqref{eq:ch05-phan-ra} của chương trước thì vector
    ngữ cảnh không có chỉ số nào. Viết ra chính xác chỗ nào trong kiến trúc
    làm cho chỉ số ấy xuất hiện, rồi trả lời: nếu ép mọi \(\alpha_{ij}\) bằng
    \(1/T\) và không cho chúng học, mô hình còn lại là mô hình gì, và nó khác
    mô hình chương~\ref{ch:encoder-decoder} ở chỗ nào?

  \item Bài viết rằng mô hình này rút về đúng RNN Encoder-Decoder của Cho
    2014a nếu cố định \(\vect{c}_i\) bằng \(\overrightarrow{\vect{h}}_{T}\).
    Kiểm câu đó bằng cách đặt giá trị ấy vào
    \eqref{eq:ch06-tong-co-trong-so} và nói bộ \(\alpha_{ij}\) nào cho ra nó.
    Rồi nói mô hình thu được khác mô hình chương~\ref{ch:encoder-decoder} ở
    chỗ nào, vì hai cái không giống nhau.

  \item Trong \eqref{eq:ch06-duong-moi} có hai số hạng. Giải thích số hạng thứ
    hai xuất hiện vì lý do gì, rồi nói vì sao mục~\ref{sec:ch06-duong-gradient}
    chỉ dùng số hạng thứ nhất để lập luận. Số hạng thứ hai lớn hay nhỏ thì lập
    luận ấy hỏng?

  \item Mục~\ref{sec:ch06-duong-gradient} đo tại lúc khởi tạo, khi mọi
    \(\alpha_{ij}\) bằng \(1/T\). Dự đoán cột \code{attention min/max} của
    bảng ấy sẽ đi về đâu sau khi huấn luyện xong, rồi nói dự đoán của bạn có
    mâu thuẫn với kết luận của mục đó không, và vì sao.

  \item Bảng ở mục~\ref{sec:ch06-be-rong} có mô hình attention đạt 1.0000 ở
    \(d_h = 32\) nhưng chỉ 0.0600 ở \(d_h = 16\). Nêu hai cơ chế khác nhau có
    thể gây ra chỗ đứt gãy sắc như thế, rồi với mỗi cơ chế nói một phép đo
    phân biệt được nó với cơ chế kia.

  \item Ở mục~\ref{sec:ch06-hai-chieu}, encoder một chiều ở \(d_h = 42\) đạt
    1.0000 còn \tn{encoder hai chiều}{bidirectional encoder} ở \(d_h = 32\)
    đạt 0.9967. Viết ra kết luận
    mà bảng ấy \emph{cho phép} rút ra, rồi viết ra một kết luận nghe rất giống
    nhưng bảng không cho phép, và nói chỗ khác nhau giữa hai câu.

  \item Hàng \code{con} trong ma trận ở mục~\ref{sec:ch06-mo-ma-tran} tản đều
    trên bốn cột, còn hàng \code{cái} dồn 0.95 vào một cột. Cả hai đều là
    \tn{loại từ}{classifier}. Nêu một giả thuyết giải thích chênh lệch ấy mà
    \emph{không} viện tới vị trí đầu câu, rồi nói phép đo nào tách giả thuyết
    của bạn khỏi giả thuyết vị trí mà mục ấy nêu.

  \item Bài Cho 2014b kết luận rằng chỗ hỏng có thể nằm ở decoder, còn bài
    Bahdanau sửa ở phía encoder. Dựng lại lập luận của bài 2014b từ dữ liệu họ
    có, rồi chỉ ra chính xác giả định nào trong lập luận ấy là giả định mà bài
    2015 phá vỡ.
\end{bt}

\tierapply

Trạng thái xuất phát cho cả năm bài: clone repo companion, đứng ở thư mục gốc
của nó, và checkout đúng tag của chương này.

\begin{minted}{console}
$ git clone https://github.com/Giang-Dang/rnn-to-transformer-lab.git
$ cd rnn-to-transformer-lab
$ git checkout ch06
$ conda env create -f environment.yml
$ conda activate rnn-to-transformer-lab
$ python verify.py --only ch06
\end{minted}

Dòng cuối phải in ra \code{verify: ok}. Nếu không, đừng làm tiếp: mọi con số
bạn đo sau đó sẽ không so được với sách. Cả năm bài dưới đây đều huấn luyện lại
ít nhất một mô hình, nên tính trước vài chục giây cho mỗi lần chạy.

\begin{bt}
  \item Bỏ mask đi. Trong \code{AdditiveAttention.weights}, xóa dòng
    \code{masked\_fill} rồi chạy \code{python -m pytest -q tests/test\_ch06.py}
    và nói test nào hỏng. Sau đó chạy lại \code{ch06\_width.py} ở riêng
    \(d_h = 32\) và ghi lại độ chính xác. Chỗ mất là bao nhiêu, và nó rơi vào
    cột \code{short} hay cột \code{long}? Giải thích chiều bạn thấy bằng cấu
    tạo của batch.

  \item \code{ch06\_gradient.py} lấy mất mát ở bước đích 5.
    Đổi \code{STEP} thành 0 rồi chạy lại. Cột \code{attention min/max} tụt
    xuống rõ rệt. Đọc \code{AttentionSeq2Seq.encode} rồi giải thích vì sao
    riêng bước 0 khác các bước khác, và nói vì sao script chọn 5 chứ không
    chọn 0.

  \item Dựng lại phép đối chứng tham số ở mục~\ref{sec:ch06-hai-chieu} theo
    chiều ngược. Thay vì nới encoder một chiều lên cho bằng, hãy bóp encoder
    hai chiều xuống: tìm \(d_h\) mà mô hình hai chiều có số tham số gần 27365
    nhất, thêm dòng ấy vào \code{RUNS} của \code{ch06\_encoder.py}, và chạy
    lại. Kết quả ủng hộ hay bác kết luận của mục đó?

  \item \code{ch06\_alignment.py} huấn luyện trên câu nguồn không đảo. Đổi
    \code{reverse\_source} thành \code{True}, chạy lại, và đọc ma trận in ra.
    Tỷ lệ bắt chéo đổi thế nào? Cẩn thận khi đọc: chỉ số cột bây giờ đếm trên
    câu đã đảo, nên hãy nói rõ \enquote{bắt chéo} còn nghĩa là gì trong hệ tọa
    độ mới trước khi so con số với 0.9744.

  \item Corpus hiện sinh câu tối đa hai mệnh đề. Cho nó sinh câu ba mệnh đề
    thay vì hai: tìm chỗ gọi hàm chia tập trong \code{ch06\_width.py} và đặt
    tham số \code{max\_clauses} của lời gọi ấy bằng 3. Giữ nguyên mọi thứ khác
    rồi chạy lại cả bảng. Ngưỡng \(d_h\) mà mô hình attention bắt đầu giải được
    dịch đi đâu, và ngưỡng của mô hình vector cố định dịch đi đâu? Hai ngưỡng
    ấy dịch cùng một lượng không, và chuyện đó ủng hộ hay bác cách đọc ở
    mục~\ref{sec:ch06-be-rong}? Ngân sách của script là 360 giây; nói rõ bạn
    đã bỏ dòng nào nếu phải bỏ.
\end{bt}

\tierextend

\begin{bt}
  \item Mục~\ref{sec:ch06-hai-chieu} kết luận rằng corpus này không cho thấy
    được chiều ngược dùng để làm gì, vì văn phạm của nó không có nhập nhằng.
    Sửa chuyện đó. Thêm vào \code{toy\_corpus.py} một chỗ nhập nhằng thật, tức
    một từ tiếng Anh mà bản dịch tiếng Việt của nó phụ thuộc vào từ đứng
    \emph{sau} nó, rồi chạy lại cả bảng encoder kèm dòng đối chứng tham số.
    Nếu chiều ngược vẫn không đáng gì, hãy nói xem chỗ nhập nhằng bạn thêm vào
    đã đủ khó chưa, và đo bằng cách nào.

  \item Mục~\ref{sec:ch06-hai-chieu} để lại một quan sát chưa giải thích được:
    với encoder một chiều, đảo câu nguồn làm độ chính xác trên câu dài tụt từ
    0.9079 xuống 0.7434, còn với encoder hai chiều thì phép đảo gần như vô
    hại. Tìm cơ chế. Gợi ý về chỗ nên nhìn trước: \(\vect{s}_0\) lấy từ đâu
    trong mỗi cấu hình, và annotation của từ nguồn đầu tiên đã đọc được những
    gì. Kết luận của bạn phải kèm một phép đo phân biệt nó với ít nhất một
    cách giải thích khác.

  \item Bài Bahdanau gọi chi phí \(T_x \times T_y\) là một \enquote{drawback}
    rồi gạt đi vì câu dịch chỉ dài 15 tới 40 từ. Biến câu ấy thành một phép đo
    trên repo: dựng một thí nghiệm đo thời gian một bước decoder theo \(T_x\),
    tách phần chi phí của attention khỏi phần chi phí của mạng hồi quy. Ở
    \(T_x\) bằng bao nhiêu thì phần attention vượt phần hồi quy? Rồi ngoại suy:
    nếu bỏ hẳn mạng hồi quy đi như chương~\ref{ch:transformer} sẽ làm, con số
    ấy đổi thế nào?

  \item Đây là chỗ còn mở thật sự. Luong và cộng sự đo được rằng alignment
    error rate và điểm dịch không đi cùng nhau, và
    mục~\ref{sec:ch06-mo-ma-tran} lấy đó làm lý do không dùng ma trận đồng
    chỉnh làm thước đo chất lượng. Nhưng tỷ lệ bắt chéo 0.9744 rõ ràng đang đo
    một cái gì đó thật. Định nghĩa một đại lượng trên
    \tn{ma trận đồng chỉnh}{alignment matrix} mà
    bạn tin là \emph{có} tương quan với chất lượng dịch trên corpus này, kiểm
    nó bằng cách huấn luyện vài mô hình ở các bề rộng khác nhau, rồi nói nó
    hỏng ở đâu. Biết trước rằng chưa có đại lượng nào thuộc loại này được chấp
    nhận rộng rãi, và rằng lý do có thể là không tồn tại đại lượng nào như thế.
\end{bt}
